% !TEX program = pdflatex
\documentclass[11pt,a4paper]{article}

% =========================================================================
% ACTAS DE SPRINT
% Proyecto: CasaScore (AlquilaPuno)
% Gestión de Proyectos TICS (EST315) -- Semestre 2026-II -- Team Sobrevivientes
% =========================================================================

% --- PAQUETES DE CONFIGURACIÓN ---
\usepackage[utf8]{inputenc}
\usepackage[T1]{fontenc}
\usepackage[spanish,es-tabla]{babel}
% Compatibilidad con babel-spanish reciente (flechas de TikZ y signo de porcentaje)
\AtBeginDocument{\shorthandoff{<>}\spanishplainpercent}
\usepackage{geometry}
\geometry{
    top=3.6cm,
    bottom=2.3cm,
    left=2.2cm,
    right=2.2cm,
    % PARTE MODIFICABLE: headheight debe ser mayor que la altura del logo más alto (2.56cm = 73pt).
    headheight=80pt,
    headsep=15pt,
    footskip=20pt
}

\usepackage{graphicx}
\usepackage{fancyhdr}
\usepackage[table]{xcolor}
\usepackage{titlesec}
\usepackage{booktabs}
\usepackage{tabularx}
\usepackage{array}
\usepackage{amsmath}
\usepackage{enumitem}
\usepackage{microtype}
\usepackage{amssymb}
% --- RECURSOS VISUALES: GRÁFICOS Y CÓDIGO DE COLORES ---
\usepackage{tikz}
\usetikzlibrary{positioning,calc}
\definecolor{GrisFila}{RGB}{246, 248, 251}
\definecolor{CuadCerca}{RGB}{150, 30, 30}
\definecolor{CuadSatisf}{RGB}{180, 130, 20}
\definecolor{CuadInfo}{RGB}{10, 45, 90}
\definecolor{CuadMonit}{RGB}{110, 115, 125}
\newcommand{\etiqueta}[2]{{\setlength{\fboxsep}{2pt}\colorbox{#1}{\textcolor{white}{\scriptsize\bfseries #2}}}}
\newcommand{\cerca}{\etiqueta{CuadCerca}{Gestionar de cerca}}
\newcommand{\satisf}{\etiqueta{CuadSatisf}{Mantener satisfecho}}
\newcommand{\informado}{\etiqueta{CuadInfo}{Mantener informado}}
\newcommand{\monit}{\etiqueta{CuadMonit}{Monitorear}}
\usepackage{needspace}
\usepackage{longtable}
\definecolor{Rama1}{RGB}{10, 45, 90}
\definecolor{Rama2}{RGB}{40, 95, 150}
\definecolor{Rama3}{RGB}{180, 130, 20}
\definecolor{Rama4}{RGB}{30, 110, 60}
\definecolor{Rama5}{RGB}{110, 60, 140}
\definecolor{RaciR}{RGB}{150, 30, 30}
\definecolor{RaciA}{RGB}{10, 45, 90}
\definecolor{RaciC}{RGB}{180, 130, 20}
\definecolor{RaciI}{RGB}{140, 145, 155}
\definecolor{NivAlto}{RGB}{30, 110, 60}
\definecolor{NivMedio}{RGB}{180, 130, 20}
\definecolor{NivBajo}{RGB}{150, 30, 30}
\usetikzlibrary{patterns}
\definecolor{SevBaja}{RGB}{30, 110, 60}
\definecolor{SevMedia}{RGB}{180, 130, 20}
\definecolor{SevAlta}{RGB}{150, 30, 30}
\usetikzlibrary{arrows.meta,shapes.geometric,fit,backgrounds}
\usepackage{listings}
\lstset{basicstyle=\ttfamily\scriptsize, frame=single, rulecolor=\color{AzulUNA!40}, backgroundcolor=\color{GrisFila}, columns=fullflexible, keepspaces=true, breaklines=true, xleftmargin=4pt, xrightmargin=4pt, literate={á}{{\'a}}1 {é}{{\'e}}1 {í}{{\'i}}1 {ó}{{\'o}}1 {ú}{{\'u}}1 {ñ}{{\~n}}1 {Á}{{\'A}}1}

\usepackage{hyperref}

% Logos: funciona a 2 o 3 niveles de profundidad dentro de docs/latex/
\graphicspath{{../../assets/img/}{../../../assets/img/}}

% --- PALETA INSTITUCIONAL FIEI / UNA PUNO ---
\definecolor{AzulUNA}{RGB}{10, 45, 90}
\definecolor{DoradoFIEI}{RGB}{180, 130, 20}
\definecolor{GrisOscuro}{RGB}{50, 50, 50}
\definecolor{AzulClaro}{RGB}{228, 236, 246}

% --- CONFIGURACIÓN DE LISTAS ---
\setlist{topsep=2pt, itemsep=2pt, parsep=1pt}

% --- COLUMNAS DE TABLA ---
\newcolumntype{L}{>{\raggedright\arraybackslash}X}
\newcolumntype{P}[1]{>{\raggedright\arraybackslash}p{#1}}
\newcolumntype{C}[1]{>{\centering\arraybackslash}p{#1}}
\newcolumntype{Q}[1]{>{\raggedleft\arraybackslash}p{#1}}
\newcolumntype{R}{>{\raggedleft\arraybackslash}X}
\renewcommand{\arraystretch}{1.25}

% --- FORMATO DE ENLACES ---
\hypersetup{
    colorlinks=true,
    linkcolor=AzulUNA,
    urlcolor=AzulUNA,
    citecolor=AzulUNA,
    pdftitle={Actas de Sprint -- CasaScore (AlquilaPuno)},
    pdfauthor={Team Sobrevivientes}
}

% --- FORMATO DE TÍTULOS ---
\titleformat{\section}{\color{AzulUNA}\normalfont\large\bfseries}{\thesection.}{0.7em}{}[\color{AzulUNA}\titlerule]
\titleformat{\subsection}{\color{AzulUNA}\normalfont\normalsize\bfseries}{\thesubsection.}{0.5em}{}
\titlespacing*{\section}{0pt}{9pt}{4pt}
\titlespacing*{\subsection}{0pt}{6pt}{2pt}

% --- ENCABEZADO Y PIE DE PÁGINA (fancyhdr) ---
\pagestyle{fancy}
\fancyhf{}

% ===== TAMAÑO DE LOS LOGOS (PARTE MODIFICABLE) =====
% logo_fiei.png  : height=1.45cm (lado izquierdo)
% logo_equipo.png: height=2.56cm (lado derecho; 25 % más grande porque la imagen tiene más borde blanco)
% Para cambiar el tamaño, modifique solo el valor height=... de cada \includegraphics.
% Si el logo más alto supera 2.6cm, aumente también headheight y top en la configuración de página.
\fancyhead[L]{%
    \includegraphics[height=1.45cm,keepaspectratio]{logo_fiei.png}%
}
\fancyhead[C]{%
    \scriptsize \textbf{\textcolor{AzulUNA}{UNIVERSIDAD NACIONAL DEL ALTIPLANO DE PUNO}}\\[2pt]
    \tiny \textcolor{GrisOscuro}{FACULTAD DE INGENIERÍA ESTADÍSTICA E INFORMÁTICA}\\[2pt]
    \scriptsize \textbf{\textcolor{AzulUNA}{Gestión de Proyectos TICS (EST315) \textbar\ Team Sobrevivientes}}\\[2pt]
    \footnotesize \textbf{\textcolor{DoradoFIEI}{ACTAS DE SPRINT}}%
}
\fancyhead[R]{%
    \includegraphics[height=2.56cm,keepaspectratio]{logo_equipo.png}%
}

\renewcommand{\headrulewidth}{0pt}
\renewcommand{\headrule}{}

\fancyfoot[L]{\scriptsize \textcolor{GrisOscuro}{CasaScore (AlquilaPuno)}}
\fancyfoot[C]{\small \thepage}
\fancyfoot[R]{\scriptsize \textcolor{GrisOscuro}{Puno, Perú -- Semestre 2026-II}}
\renewcommand{\footrulewidth}{0.4pt}

\begin{document}

\begin{center}
    {\Large \textbf{\textcolor{AzulUNA}{ACTAS DE SPRINT}}}\\[3pt]
    {\normalsize \textbf{CasaScore (AlquilaPuno)}}\\[4pt]
    \footnotesize \textbf{Periodo:} semanas 4 a 8 (14/09 -- 16/10/2026) \quad \textbar\quad \textbf{Versión:} 1.0 \quad \textbar\quad \textbf{Fase:} Implementación
\end{center}

\vspace{0.1cm}

\section{Datos Generales}
\noindent\begin{tabularx}{\textwidth}{@{}P{4.2cm} L@{}}
\toprule
\textbf{Proyecto:} & CasaScore (aplicación AlquilaPuno) \\
\textbf{Responsable:} & Viviana Mary Quisbert Quispe -- Gestora de Comunicación, Riesgos y Documentación \\
\textbf{Proceso PM4DEV:} & Gestión de la información (seguimiento del trabajo) \\
\textbf{Reunión:} & Revisión del sprint, viernes de 20:00 a 21:00 por Google Meet \\
\textbf{Asistentes:} & Los cuatro integrantes del equipo en todas las sesiones \\
\bottomrule
\end{tabularx}

\section{Cómo Leer las Actas}
Cada sprint dura una semana. En la tabla de avance, la línea azul marca el avance planificado según la línea base del cronograma y la barra muestra el avance real: \textcolor{Rama4}{\textbf{verde}} si está al día o adelantado, \textcolor{SevMedia}{\textbf{ámbar}} si está hasta 20 puntos por debajo y \textcolor{SevAlta}{\textbf{rojo}} si el atraso es mayor.

\vspace{0.15cm}
{\small
\rowcolors{2}{white}{GrisFila}
\noindent\begin{tabularx}{\textwidth}{@{}C{1.3cm} C{1.0cm} P{2.4cm} L C{1.4cm}@{}}
\toprule
\rowcolor{AzulClaro}
\textbf{Sprint} & \textbf{Sem.} & \textbf{Fecha} & \textbf{Objetivo} & \textbf{SPI} \\
\midrule


1 & 4 & 18/09/2026 & Terminar la planificación base y avanzar el motor CASA SCORE. & 0.78 \\

2 & 5 & 25/09/2026 & Cerrar el motor CASA SCORE y el modelo de datos; registrar la decisión sobre Horus Sentinel. & 0.65 \\

3 & 6 & 02/10/2026 & Terminar la publicación de cuartos e iniciar el levantamiento de cuartos reales. & 0.56 \\

4 & 7 & 09/10/2026 & Entregar el prototipo funcional con ranking explicado (MVP de la Unidad I). & 0.62 \\

5 & 8 & 16/10/2026 & Sustentar la Unidad I y cerrar el ranking y el detalle. & 0.66 \\

\bottomrule
\end{tabularx}
}


\Needspace{20\baselineskip}
\section{Acta del Sprint 1 -- Semana 4}
\noindent\textbf{Fecha:} viernes 18/09/2026 \quad \textbar\quad \textbf{Objetivo del sprint:} Terminar la planificación base y avanzar el motor CASA SCORE.

\vspace{0.1cm}
\noindent\textbf{Avance real}

\vspace{0.05cm}
{\footnotesize
\setlength{\tabcolsep}{3pt}
\rowcolors{2}{white}{GrisFila}
\noindent\begin{tabularx}{\textwidth}{@{}C{0.7cm} P{4.2cm} C{2.5cm} C{0.85cm} C{0.85cm} L@{}}
\toprule
\rowcolor{AzulClaro}
\textbf{Cód.} & \textbf{Actividad} & \textbf{Avance} & \textbf{Plan} & \textbf{Real} & \textbf{Resultado} \\
\midrule
1.3 & Planificación base (interesados, EDT, equipo, cronograma y presupuesto) & \tikz[baseline=-0.6ex]{\fill[gray!20] (0,-0.09) rectangle (2.20,0.09); \draw[AzulUNA, line width=1.3pt] (2.090,-0.14) -- (2.090,0.14); \fill[SevAlta] (0,-0.06) rectangle (0.660,0.06);} & 95\,\% & \textbf{30\,\%} & Registro de interesados terminado; EDT y plan del equipo en borrador. \\
3.1 & Motor CASA SCORE & \tikz[baseline=-0.6ex]{\fill[gray!20] (0,-0.09) rectangle (2.20,0.09); \draw[AzulUNA, line width=1.3pt] (0.990,-0.14) -- (0.990,0.14); \fill[Rama4] (0,-0.06) rectangle (1.100,0.06);} & 45\,\% & \textbf{50\,\%} & Normalización de precio, cercanía, seguridad y servicios implementada. \\
2.2 & Modelo de datos y arquitectura & \tikz[baseline=-0.6ex]{\fill[gray!20] (0,-0.09) rectangle (2.20,0.09); \draw[AzulUNA, line width=1.3pt] (2.090,-0.14) -- (2.090,0.14); \fill[SevAlta] (0,-0.06) rectangle (0.660,0.06);} & 95\,\% & \textbf{30\,\%} & Tabla local de cuartos definida. \\
2.3 & Diseño de pantallas & \tikz[baseline=-0.6ex]{\fill[gray!20] (0,-0.09) rectangle (2.20,0.09); \draw[AzulUNA, line width=1.3pt] (2.090,-0.14) -- (2.090,0.14); \fill[SevAlta] (0,-0.06) rectangle (0.440,0.06);} & 95\,\% & \textbf{20\,\%} & Bocetos de ranking y detalle. \\
\bottomrule
\end{tabularx}
}

\vspace{0.1cm}
\noindent\textbf{Desviaciones}
\begin{itemize}[leftmargin=*, itemsep=0pt]
    \item La planificación base avanza menos de lo previsto porque el equipo elabora en paralelo los documentos de Horus Sentinel.
    \item Los primeros cálculos del cronograma y del presupuesto de Horus Sentinel indican que no cabe en el semestre.
\end{itemize}
\noindent\textbf{Acuerdos}
\begin{enumerate}[leftmargin=*, itemsep=0pt]
    \item CasaScore tiene prioridad sobre Horus Sentinel cuando compitan por tiempo.
    \item La continuidad de Horus Sentinel se decide el martes 22/09 con su cronograma y presupuesto detallados.
\end{enumerate}
\noindent\textbf{Compromisos}

\vspace{0.05cm}
{\footnotesize
\rowcolors{2}{white}{GrisFila}
\noindent\begin{tabularx}{\textwidth}{@{}L P{3.6cm} C{1.6cm}@{}}
\toprule
\rowcolor{AzulClaro}
\textbf{Compromiso} & \textbf{Responsable} & \textbf{Fecha} \\
\midrule
Terminar el cronograma y el presupuesto de ambos proyectos & Ruth K. Apaza Solis & 22/09 \\
Completar el motor CASA SCORE con bonificaciones y penalizaciones & Eliseo Tarqui Ajahuana & 25/09 \\
Preparar el análisis de continuidad de Horus Sentinel & Clyde N. Paricahua Pari & 22/09 \\
\bottomrule
\end{tabularx}
}

\Needspace{20\baselineskip}
\section{Acta del Sprint 2 -- Semana 5}
\noindent\textbf{Fecha:} viernes 25/09/2026 \quad \textbar\quad \textbf{Objetivo del sprint:} Cerrar el motor CASA SCORE y el modelo de datos; registrar la decisión sobre Horus Sentinel.

\vspace{0.1cm}
\noindent\textbf{Avance real}

\vspace{0.05cm}
{\footnotesize
\setlength{\tabcolsep}{3pt}
\rowcolors{2}{white}{GrisFila}
\noindent\begin{tabularx}{\textwidth}{@{}C{0.7cm} P{4.2cm} C{2.5cm} C{0.85cm} C{0.85cm} L@{}}
\toprule
\rowcolor{AzulClaro}
\textbf{Cód.} & \textbf{Actividad} & \textbf{Avance} & \textbf{Plan} & \textbf{Real} & \textbf{Resultado} \\
\midrule
1.3 & Planificación base & \tikz[baseline=-0.6ex]{\fill[gray!20] (0,-0.09) rectangle (2.20,0.09); \draw[AzulUNA, line width=1.3pt] (2.200,-0.14) -- (2.200,0.14); \fill[SevAlta] (0,-0.06) rectangle (0.880,0.06);} & 100\,\% & \textbf{40\,\%} & EDT, plan del equipo, cronograma y presupuesto en revisión. \\
1.4 & Planes de soporte (riesgos, calidad, comunicaciones, contratos y arquitectura) & \tikz[baseline=-0.6ex]{\fill[gray!20] (0,-0.09) rectangle (2.20,0.09); \draw[AzulUNA, line width=1.3pt] (1.760,-0.14) -- (1.760,0.14); \fill[SevAlta] (0,-0.06) rectangle (0.110,0.06);} & 80\,\% & \textbf{5\,\%} & Borradores iniciados. \\
3.1 & Motor CASA SCORE & \tikz[baseline=-0.6ex]{\fill[gray!20] (0,-0.09) rectangle (2.20,0.09); \draw[AzulUNA, line width=1.3pt] (2.090,-0.14) -- (2.090,0.14); \fill[Rama4] (0,-0.06) rectangle (2.200,0.06);} & 95\,\% & \textbf{100\,\%} & Terminado y conectado a la base local SQLite. \\
2.2 & Modelo de datos y arquitectura & \tikz[baseline=-0.6ex]{\fill[gray!20] (0,-0.09) rectangle (2.20,0.09); \draw[AzulUNA, line width=1.3pt] (2.200,-0.14) -- (2.200,0.14); \fill[SevAlta] (0,-0.06) rectangle (1.100,0.06);} & 100\,\% & \textbf{50\,\%} & Modelo local terminado; modelo en la nube pendiente. \\
2.3 & Diseño de pantallas & \tikz[baseline=-0.6ex]{\fill[gray!20] (0,-0.09) rectangle (2.20,0.09); \draw[AzulUNA, line width=1.3pt] (2.200,-0.14) -- (2.200,0.14); \fill[SevAlta] (0,-0.06) rectangle (0.880,0.06);} & 100\,\% & \textbf{40\,\%} & Ranking y detalle definidos; publicar y perfil pendientes. \\
3.2 & Publicación de cuartos & \tikz[baseline=-0.6ex]{\fill[gray!20] (0,-0.09) rectangle (2.20,0.09); \draw[AzulUNA, line width=1.3pt] (1.870,-0.14) -- (1.870,0.14); \fill[SevAlta] (0,-0.06) rectangle (0.440,0.06);} & 85\,\% & \textbf{20\,\%} & Estructura del formulario iniciada. \\
\bottomrule
\end{tabularx}
}

\vspace{0.1cm}
\noindent\textbf{Desviaciones}
\begin{itemize}[leftmargin=*, itemsep=0pt]
    \item La gestión acumula un atraso de una semana y media: los documentos de planificación siguen en revisión.
    \item El motor de puntaje se terminó antes de lo previsto.
\end{itemize}
\noindent\textbf{Acuerdos}
\begin{enumerate}[leftmargin=*, itemsep=0pt]
    \item Se registra la exclusión de Horus Sentinel aprobada el martes 22/09 en su acta.
    \item La Analista de Planificación preparará una solicitud de cambio para evaluar más horas semanales y recuperar el atraso.
\end{enumerate}
\noindent\textbf{Compromisos}

\vspace{0.05cm}
{\footnotesize
\rowcolors{2}{white}{GrisFila}
\noindent\begin{tabularx}{\textwidth}{@{}L P{3.6cm} C{1.6cm}@{}}
\toprule
\rowcolor{AzulClaro}
\textbf{Compromiso} & \textbf{Responsable} & \textbf{Fecha} \\
\midrule
Presentar la solicitud de cambio con su análisis de impacto & Ruth K. Apaza Solis & 29/09 \\
Terminar el formulario de publicación con 4 fotos & Eliseo Tarqui Ajahuana & 02/10 \\
Visitar propietarios para levantar los primeros cuartos & Clyde N. Paricahua Pari & 02/10 \\
\bottomrule
\end{tabularx}
}

\Needspace{20\baselineskip}
\section{Acta del Sprint 3 -- Semana 6}
\noindent\textbf{Fecha:} viernes 02/10/2026 \quad \textbar\quad \textbf{Objetivo del sprint:} Terminar la publicación de cuartos e iniciar el levantamiento de cuartos reales.

\vspace{0.1cm}
\noindent\textbf{Avance real}

\vspace{0.05cm}
{\footnotesize
\setlength{\tabcolsep}{3pt}
\rowcolors{2}{white}{GrisFila}
\noindent\begin{tabularx}{\textwidth}{@{}C{0.7cm} P{4.2cm} C{2.5cm} C{0.85cm} C{0.85cm} L@{}}
\toprule
\rowcolor{AzulClaro}
\textbf{Cód.} & \textbf{Actividad} & \textbf{Avance} & \textbf{Plan} & \textbf{Real} & \textbf{Resultado} \\
\midrule
1.3 & Planificación base & \tikz[baseline=-0.6ex]{\fill[gray!20] (0,-0.09) rectangle (2.20,0.09); \draw[AzulUNA, line width=1.3pt] (2.200,-0.14) -- (2.200,0.14); \fill[SevAlta] (0,-0.06) rectangle (0.880,0.06);} & 100\,\% & \textbf{40\,\%} & Sin cambios: documentos en revisión. \\
1.4 & Planes de soporte & \tikz[baseline=-0.6ex]{\fill[gray!20] (0,-0.09) rectangle (2.20,0.09); \draw[AzulUNA, line width=1.3pt] (2.200,-0.14) -- (2.200,0.14); \fill[SevAlta] (0,-0.06) rectangle (0.220,0.06);} & 100\,\% & \textbf{10\,\%} & Borradores de riesgos y calidad. \\
3.2 & Publicación de cuartos & \tikz[baseline=-0.6ex]{\fill[gray!20] (0,-0.09) rectangle (2.20,0.09); \draw[AzulUNA, line width=1.3pt] (2.200,-0.14) -- (2.200,0.14); \fill[SevAlta] (0,-0.06) rectangle (1.320,0.06);} & 100\,\% & \textbf{60\,\%} & Formulario con 4 fotos y campos obligatorios; falta validar con propietarios. \\
4.2 & Levantamiento de cuartos & \tikz[baseline=-0.6ex]{\fill[gray!20] (0,-0.09) rectangle (2.20,0.09); \draw[AzulUNA, line width=1.3pt] (2.200,-0.14) -- (2.200,0.14); \fill[SevAlta] (0,-0.06) rectangle (0.550,0.06);} & 100\,\% & \textbf{25\,\%} & 4 de 15 cuartos registrados. \\
2.2 & Modelo de datos y arquitectura & \tikz[baseline=-0.6ex]{\fill[gray!20] (0,-0.09) rectangle (2.20,0.09); \draw[AzulUNA, line width=1.3pt] (2.200,-0.14) -- (2.200,0.14); \fill[SevAlta] (0,-0.06) rectangle (1.320,0.06);} & 100\,\% & \textbf{60\,\%} & Entidades de la nube definidas. \\
2.3 & Diseño de pantallas & \tikz[baseline=-0.6ex]{\fill[gray!20] (0,-0.09) rectangle (2.20,0.09); \draw[AzulUNA, line width=1.3pt] (2.200,-0.14) -- (2.200,0.14); \fill[SevAlta] (0,-0.06) rectangle (1.100,0.06);} & 100\,\% & \textbf{50\,\%} & Pantalla de publicar terminada. \\
\bottomrule
\end{tabularx}
}

\vspace{0.1cm}
\noindent\textbf{Desviaciones}
\begin{itemize}[leftmargin=*, itemsep=0pt]
    \item El valor ganado al corte es de S/ 2,623.86 frente a S/ 4,691.79 planificados (SPI de 0.56): el atraso equivale a dos semanas.
    \item Las horas reales registradas son menores a las planificadas por la carga de Horus Sentinel hasta la semana 5 y de otras asignaturas.
    \item El costo se mantiene controlado (CPI de 0.96).
\end{itemize}
\noindent\textbf{Acuerdos}
\begin{enumerate}[leftmargin=*, itemsep=0pt]
    \item El martes 29/09 se evaluó la solicitud de cambio para pasar de 8 a 10 horas semanales y fue rechazada por unanimidad.
    \item Se aprueba un plan de recuperación sin horas adicionales: construir el ranking en paralelo con la documentación y levantar cuartos los sábados.
    \item La revisión y firma de los documentos de planificación se hará en las semanas 9 y 10, cuando el desarrollo se detiene por los exámenes.
\end{enumerate}
\noindent\textbf{Compromisos}

\vspace{0.05cm}
{\footnotesize
\rowcolors{2}{white}{GrisFila}
\noindent\begin{tabularx}{\textwidth}{@{}L P{3.6cm} C{1.6cm}@{}}
\toprule
\rowcolor{AzulClaro}
\textbf{Compromiso} & \textbf{Responsable} & \textbf{Fecha} \\
\midrule
Construir el ranking y el detalle del cuarto & Eliseo Tarqui Ajahuana & 09/10 \\
Ejecutar las pruebas del motor con 5 cuartos de prueba & Viviana M. Quisbert Quispe & 09/10 \\
Registrar 4 cuartos más & Clyde N. Paricahua Pari & 09/10 \\
\bottomrule
\end{tabularx}
}

\Needspace{20\baselineskip}
\section{Acta del Sprint 4 -- Semana 7}
\noindent\textbf{Fecha:} viernes 09/10/2026 \quad \textbar\quad \textbf{Objetivo del sprint:} Entregar el prototipo funcional con ranking explicado (MVP de la Unidad I).

\vspace{0.1cm}
\noindent\textbf{Avance real}

\vspace{0.05cm}
{\footnotesize
\setlength{\tabcolsep}{3pt}
\rowcolors{2}{white}{GrisFila}
\noindent\begin{tabularx}{\textwidth}{@{}C{0.7cm} P{4.2cm} C{2.5cm} C{0.85cm} C{0.85cm} L@{}}
\toprule
\rowcolor{AzulClaro}
\textbf{Cód.} & \textbf{Actividad} & \textbf{Avance} & \textbf{Plan} & \textbf{Real} & \textbf{Resultado} \\
\midrule
3.2 & Publicación de cuartos & \tikz[baseline=-0.6ex]{\fill[gray!20] (0,-0.09) rectangle (2.20,0.09); \draw[AzulUNA, line width=1.3pt] (2.200,-0.14) -- (2.200,0.14); \fill[Rama4] (0,-0.06) rectangle (2.200,0.06);} & 100\,\% & \textbf{100\,\%} & Terminado. \\
3.3 & Ranking y detalle (MVP de la Unidad I) & \tikz[baseline=-0.6ex]{\fill[gray!20] (0,-0.09) rectangle (2.20,0.09); \draw[AzulUNA, line width=1.3pt] (0.990,-0.14) -- (0.990,0.14); \fill[Rama4] (0,-0.06) rectangle (1.760,0.06);} & 45\,\% & \textbf{80\,\%} & Ranking con 4 perfiles de prioridades y desglose del puntaje. \\
4.1 & Pruebas del motor & \tikz[baseline=-0.6ex]{\fill[gray!20] (0,-0.09) rectangle (2.20,0.09); \draw[AzulUNA, line width=1.3pt] (2.200,-0.14) -- (2.200,0.14); \fill[SevAlta] (0,-0.06) rectangle (1.540,0.06);} & 100\,\% & \textbf{70\,\%} & Casos ejecutados con 5 cuartos de prueba; pruebas automáticas pendientes. \\
4.2 & Levantamiento de cuartos & \tikz[baseline=-0.6ex]{\fill[gray!20] (0,-0.09) rectangle (2.20,0.09); \draw[AzulUNA, line width=1.3pt] (2.200,-0.14) -- (2.200,0.14); \fill[SevAlta] (0,-0.06) rectangle (0.880,0.06);} & 100\,\% & \textbf{40\,\%} & 6 de 15 cuartos registrados. \\
2.2 & Modelo de datos y arquitectura & \tikz[baseline=-0.6ex]{\fill[gray!20] (0,-0.09) rectangle (2.20,0.09); \draw[AzulUNA, line width=1.3pt] (2.200,-0.14) -- (2.200,0.14); \fill[SevMedia] (0,-0.06) rectangle (1.760,0.06);} & 100\,\% & \textbf{80\,\%} & Modelo entidad--relación de la nube terminado. \\
1.3 & Planificación base & \tikz[baseline=-0.6ex]{\fill[gray!20] (0,-0.09) rectangle (2.20,0.09); \draw[AzulUNA, line width=1.3pt] (2.200,-0.14) -- (2.200,0.14); \fill[SevAlta] (0,-0.06) rectangle (1.100,0.06);} & 100\,\% & \textbf{50\,\%} & Cronograma y presupuesto revisados. \\
\bottomrule
\end{tabularx}
}

\vspace{0.1cm}
\noindent\textbf{Desviaciones}
\begin{itemize}[leftmargin=*, itemsep=0pt]
    \item El prototipo se entregó el 09/10 con el informe de validación, tres días antes de la línea base del MVP de la Unidad I.
    \item La gestión mantiene el atraso de dos semanas.
    \item El puntaje llega al tope de 100 en el mejor cuarto en todos los perfiles, porque las bonificaciones superan el margen disponible.
\end{itemize}
\noindent\textbf{Acuerdos}
\begin{enumerate}[leftmargin=*, itemsep=0pt]
    \item El prototipo queda como base del MVP de la Unidad I.
    \item El tope del puntaje se revisará en el ajuste del motor del MVP final.
\end{enumerate}
\noindent\textbf{Compromisos}

\vspace{0.05cm}
{\footnotesize
\rowcolors{2}{white}{GrisFila}
\noindent\begin{tabularx}{\textwidth}{@{}L P{3.6cm} C{1.6cm}@{}}
\toprule
\rowcolor{AzulClaro}
\textbf{Compromiso} & \textbf{Responsable} & \textbf{Fecha} \\
\midrule
Preparar la sustentación de la Unidad I & Clyde N. Paricahua Pari & 16/10 \\
Completar el levantamiento hasta 9 cuartos & Clyde N. Paricahua Pari & 16/10 \\
Implementar las pruebas automáticas con Jest & Eliseo Tarqui Ajahuana & 16/10 \\
\bottomrule
\end{tabularx}
}

\Needspace{20\baselineskip}
\section{Acta del Sprint 5 -- Semana 8}
\noindent\textbf{Fecha:} viernes 16/10/2026 \quad \textbar\quad \textbf{Objetivo del sprint:} Sustentar la Unidad I y cerrar el ranking y el detalle.

\vspace{0.1cm}
\noindent\textbf{Avance real}

\vspace{0.05cm}
{\footnotesize
\setlength{\tabcolsep}{3pt}
\rowcolors{2}{white}{GrisFila}
\noindent\begin{tabularx}{\textwidth}{@{}C{0.7cm} P{4.2cm} C{2.5cm} C{0.85cm} C{0.85cm} L@{}}
\toprule
\rowcolor{AzulClaro}
\textbf{Cód.} & \textbf{Actividad} & \textbf{Avance} & \textbf{Plan} & \textbf{Real} & \textbf{Resultado} \\
\midrule
3.3 & Ranking y detalle (MVP de la Unidad I) & \tikz[baseline=-0.6ex]{\fill[gray!20] (0,-0.09) rectangle (2.20,0.09); \draw[AzulUNA, line width=1.3pt] (2.200,-0.14) -- (2.200,0.14); \fill[Rama4] (0,-0.06) rectangle (2.200,0.06);} & 100\,\% & \textbf{100\,\%} & Terminado y sustentado. \\
4.1 & Pruebas del motor & \tikz[baseline=-0.6ex]{\fill[gray!20] (0,-0.09) rectangle (2.20,0.09); \draw[AzulUNA, line width=1.3pt] (2.200,-0.14) -- (2.200,0.14); \fill[SevMedia] (0,-0.06) rectangle (1.760,0.06);} & 100\,\% & \textbf{80\,\%} & Pruebas automáticas iniciadas. \\
4.2 & Levantamiento de cuartos & \tikz[baseline=-0.6ex]{\fill[gray!20] (0,-0.09) rectangle (2.20,0.09); \draw[AzulUNA, line width=1.3pt] (2.200,-0.14) -- (2.200,0.14); \fill[SevAlta] (0,-0.06) rectangle (1.320,0.06);} & 100\,\% & \textbf{60\,\%} & 9 de 15 cuartos registrados. \\
1.3 & Planificación base & \tikz[baseline=-0.6ex]{\fill[gray!20] (0,-0.09) rectangle (2.20,0.09); \draw[AzulUNA, line width=1.3pt] (2.200,-0.14) -- (2.200,0.14); \fill[SevAlta] (0,-0.06) rectangle (1.100,0.06);} & 100\,\% & \textbf{50\,\%} & Pendiente la firma de los documentos. \\
1.4 & Planes de soporte & \tikz[baseline=-0.6ex]{\fill[gray!20] (0,-0.09) rectangle (2.20,0.09); \draw[AzulUNA, line width=1.3pt] (2.200,-0.14) -- (2.200,0.14); \fill[SevAlta] (0,-0.06) rectangle (0.440,0.06);} & 100\,\% & \textbf{20\,\%} & Borradores de los cinco planes. \\
1.6 & Entregas de la Unidad II & \tikz[baseline=-0.6ex]{\fill[gray!20] (0,-0.09) rectangle (2.20,0.09); \draw[AzulUNA, line width=1.3pt] (0.110,-0.14) -- (0.110,0.14); \fill[SevMedia] (0,-0.06) rectangle (0.000,0.06);} & 5\,\% & \textbf{0\,\%} & Inicia la próxima semana. \\
\bottomrule
\end{tabularx}
}

\vspace{0.1cm}
\noindent\textbf{Desviaciones}
\begin{itemize}[leftmargin=*, itemsep=0pt]
    \item El software cumple el hito de la Unidad I; la gestión sigue dos semanas atrasada.
    \item La migración a Supabase, prevista desde el 13/10, se inicia en la semana 10 por la semana de exámenes; se vigila el disparador del riesgo de retraso en la nube.
\end{itemize}
\noindent\textbf{Acuerdos}
\begin{enumerate}[leftmargin=*, itemsep=0pt]
    \item En la semana 9 se elaboran la guía de marcos de referencia y la revisión de los documentos de planificación.
    \item En la semana 10 se presenta el Plan de Dirección del Proyecto y se inicia la migración a Supabase.
    \item El levantamiento de cuartos se completa hasta la semana 12, antes del piloto.
\end{enumerate}
\noindent\textbf{Compromisos}

\vspace{0.05cm}
{\footnotesize
\rowcolors{2}{white}{GrisFila}
\noindent\begin{tabularx}{\textwidth}{@{}L P{3.6cm} C{1.6cm}@{}}
\toprule
\rowcolor{AzulClaro}
\textbf{Compromiso} & \textbf{Responsable} & \textbf{Fecha} \\
\midrule
Elaborar la guía de marcos de referencia, PMBOK y arquitecturas & Clyde N. Paricahua Pari & 23/10 \\
Revisar y firmar los documentos de planificación & Viviana M. Quisbert Quispe & 30/10 \\
Iniciar la migración a Supabase & Eliseo Tarqui Ajahuana & 30/10 \\
\bottomrule
\end{tabularx}
}

\Needspace{14\baselineskip}
\section{Evolución del Desempeño del Cronograma}
El SPI compara el trabajo realizado con el planificado: 1.00 significa al día y un valor menor indica atraso.

\begin{center}
\begin{tikzpicture}[x=2.3cm,y=4cm,font=\scriptsize]
\fill[Rama4!10] (3.6,0.9) rectangle (8.4,1.1); \fill[SevMedia!10] (3.6,0.7) rectangle (8.4,0.9); \fill[SevAlta!8] (3.6,0.4) rectangle (8.4,0.7);
\draw[gray!50] (3.6,0.4) -- (8.4,0.4);
\draw[gray!25] (3.6,0.4) -- (8.4,0.4); \node[left, font=\tiny] at (3.6,0.4) {0.40};
\draw[gray!25] (3.6,0.6) -- (8.4,0.6); \node[left, font=\tiny] at (3.6,0.6) {0.60};
\draw[gray!25] (3.6,0.8) -- (8.4,0.8); \node[left, font=\tiny] at (3.6,0.8) {0.80};
\draw[gray!25] (3.6,1.0) -- (8.4,1.0); \node[left, font=\tiny] at (3.6,1.0) {1.00};
\node[below, font=\tiny] at (4,0.4) {Semana 4};
\node[below, font=\tiny] at (5,0.4) {Semana 5};
\node[below, font=\tiny] at (6,0.4) {Semana 6};
\node[below, font=\tiny] at (7,0.4) {Semana 7};
\node[below, font=\tiny] at (8,0.4) {Semana 8};
\draw[AzulUNA, dashed] (3.6,1) -- (8.4,1) node[right, font=\tiny] {al día};
\draw[SevAlta, very thick] (4,0.78) -- (5,0.65) -- (6,0.56) -- (7,0.62) -- (8,0.66);
\fill[SevAlta] (4,0.78) circle (2pt); \node[above, font=\tiny\bfseries] at (4,0.78) {0.78};
\fill[SevAlta] (5,0.65) circle (2pt); \node[above, font=\tiny\bfseries] at (5,0.65) {0.65};
\fill[SevAlta] (6,0.56) circle (2pt); \node[above, font=\tiny\bfseries] at (6,0.56) {0.56};
\fill[SevAlta] (7,0.62) circle (2pt); \node[above, font=\tiny\bfseries] at (7,0.62) {0.62};
\fill[SevAlta] (8,0.66) circle (2pt); \node[above, font=\tiny\bfseries] at (8,0.66) {0.66};
\end{tikzpicture}
\end{center}
\Needspace{24\baselineskip}
\section{Firmas}
\noindent En señal de conformidad con el presente documento, los integrantes del equipo firman a continuación:

\vspace{1.5cm}
\begin{center}
\begin{tabular}{c c c}
    \parbox{6.5cm}{\centering \rule{6.0cm}{0.6pt}\\[5pt] \textbf{Clyde Neil Paricahua Pari}\\[2pt] {\footnotesize Project Manager}}
    & \hspace{0.8cm} &
    \parbox{6.5cm}{\centering \rule{6.0cm}{0.6pt}\\[5pt] \textbf{Ruth Karina Apaza Solis}\\[2pt] {\footnotesize Analista de Planificación}}
    \\[2.0cm]
    \parbox{6.5cm}{\centering \rule{6.0cm}{0.6pt}\\[5pt] \textbf{Eliseo Tarqui Ajahuana}\\[2pt] {\footnotesize Especialista Técnico / Arquitectura}}
    & \hspace{0.8cm} &
    \parbox{6.5cm}{\centering \rule{6.0cm}{0.6pt}\\[5pt] \textbf{Viviana Mary Quisbert Quispe}\\[2pt] {\footnotesize Gestora de Comunicación, Riesgos y Documentación}}
\end{tabular}
\end{center}
\end{document}
